%! Linear Algebra — Cumulative Final Exam — Final (blank)
\documentclass[10pt]{article}
\usepackage{linalg-article}
\usepackage{linalg-boxes}
\newcommand{\vv}{\mathbf{v}}
\newcommand{\ww}{\mathbf{w}}
\newcommand{\xx}{\mathbf{x}}
\newcommand{\bb}{\mathbf{b}}
\newcommand{\pp}{\mathbf{p}}
\newcommand{\ee}{\mathbf{e}}
\newcommand{\zero}{\mathbf{0}}
\newcommand{\uu}{\mathbf{u}}
\newcommand{\avec}{\mathbf{a}}
\newcommand{\relu}{\mathrm{ReLU}}
\newcommand{\T}{\mathsf{T}}

% Part-divider strip
\newcommand{\parthead}[1]{%
  \vspace{6pt}\noindent
  \begin{headlinebox}{burgundy}{\color{white}\bfseries #1}\end{headlinebox}%
  \vspace{4pt}}

\begin{document}

\pageheader{Linear Algebra --- Cumulative Final}{Final Exam}
\namedateperiod

\vspace{0.06in}
\noindent\textbf{Instructions:} Show all work. No notes unless your teacher says so. The exam has
$50$ questions in four parts --- A Vocabulary, B Multiple Choice, C Short Answer \& Computation,
D Extended Response. \hfill\textbf{Total: 100 pts.}

% ======================================================================
\parthead{Part A --- Vocabulary (16 pts)}
Write the letter of the best description next to each term. Two matching sets.

\vspace{2pt}
\noindent\textit{Set 1 --- Units 1--4.}\\[2pt]
\noindent
\begin{minipage}[t]{0.40\textwidth}
\begin{enumerate}[leftmargin=2.2em, itemsep=5pt, label=\arabic*.]
  \item Dot product \blank{1.1cm}
  \item Unit vector \blank{1.1cm}
  \item Basis \blank{1.1cm}
  \item Nullspace $N(A)$ \blank{1.1cm}
  \item Pivot \blank{1.1cm}
  \item Least squares \blank{1.1cm}
  \item Inverse matrix \blank{1.1cm}
  \item Projection \blank{1.1cm}
\end{enumerate}
\end{minipage}%
\hfill
\begin{minipage}[t]{0.57\textwidth}
\small
\begin{description}[leftmargin=1.4em, itemsep=2pt]
  \item[(A)] A vector whose length is exactly $1$.
  \item[(B)] The first nonzero entry in a row, used to clear the entries below it.
  \item[(C)] $\vv\cdot\ww=v_1w_1+v_2w_2$; it measures length and angle.
  \item[(D)] The set of all solutions to $A\xx=\zero$.
  \item[(E)] $A^{-1}$, the matrix that undoes $A$: $A^{-1}A=I$.
  \item[(F)] An independent set that spans a space --- the fewest vectors that build every vector in it.
  \item[(G)] Choosing $\hat\xx$ to make the error $\|\bb-A\xx\|^2$ as small as possible.
  \item[(H)] Sending $\bb$ to the closest point $\pp$ in a subspace.
\end{description}
\end{minipage}

\vspace{6pt}
\noindent\textit{Set 2 --- Units 5--8.}\\[2pt]
\noindent
\begin{minipage}[t]{0.40\textwidth}
\begin{enumerate}[leftmargin=2.2em, itemsep=5pt, label=\arabic*., start=9]
  \item Eigenvector \blank{1.1cm}
  \item Determinant \blank{1.1cm}
  \item Covariance matrix \blank{1.1cm}
  \item Singular value \blank{1.1cm}
  \item ReLU \blank{1.1cm}
  \item Linear transformation \blank{1.1cm}
  \item Diagonalization \blank{1.1cm}
  \item Principal component \blank{1.1cm}
\end{enumerate}
\end{minipage}%
\hfill
\begin{minipage}[t]{0.57\textwidth}
\small
\begin{description}[leftmargin=1.4em, itemsep=2pt]
  \item[(A)] A nonzero vector $\xx$ whose direction is unchanged by $A$: $A\xx=\lambda\xx$.
  \item[(B)] The symmetric matrix of variances and covariances; its eigenvectors are the principal components.
  \item[(C)] The map $\xx\mapsto A\xx$ that fixes the origin and sends grid lines to straight, evenly spaced lines.
  \item[(D)] A number whose absolute value is the area (or volume) factor; it is $0$ exactly when the matrix is singular.
  \item[(E)] The rectifier $\max(0,x)$: keep positive inputs, send negatives to $0$.
  \item[(F)] A number $\sigma\ge0$: how far $A$ stretches along a singular direction ($\sigma=\sqrt\lambda$ of $A^\T A$).
  \item[(G)] Writing $A=X\Lambda X^{-1}$ with the eigenvectors in $X$ and the eigenvalues down $\Lambda$.
  \item[(H)] A top singular direction of centered data --- the direction of greatest spread.
\end{description}
\end{minipage}

% ======================================================================
\parthead{Part B --- Multiple Choice (24 pts)}
Circle the best answer. ($2$ pts each.)

\begin{enumerate}[leftmargin=1.9em, itemsep=8pt]
  \item The length of the vector $(5,12)$ is
    \hfill (A) $17$ \quad (B) $13$ \quad (C) $60$ \quad (D) $7$
  \item Saying two nonzero vectors have dot product $0$ means they are
    \hfill (A) pointing the same way \quad (B) perpendicular \quad (C) equal \quad (D) both zero
  \item A $2\times2$ matrix whose determinant $ad-bc=0$
    \hfill (A) has a unique inverse \quad (B) is orthogonal \quad (C) is singular (no inverse) \quad (D) is symmetric
  \item For an $m\times n$ matrix of rank $r$, the dimension of the column space $C(A)$ is
    \hfill (A) $n-r$ \quad (B) $m$ \quad (C) $r$ \quad (D) $n$
  \item For a matrix with $n$ columns, $\dim C(A)+\dim N(A)$ equals
    \hfill (A) $m$ \quad (B) $n$ \quad (C) $r$ \quad (D) $0$
  \item The projection $\pp$ of $\bb$ onto $C(A)$ is the point of $C(A)$ that is
    \hfill (A) farthest from $\bb$ \quad (B) closest to $\bb$ \quad (C) equal to $\bb$ \quad (D) at the origin
  \item The quantity $|\det A|$ measures
    \hfill (A) the trace \quad (B) the area/volume scaling factor \quad (C) the rank \quad (D) the largest entry
  \item If $\det A=0$, the transformation $\xx\mapsto A\xx$
    \hfill (A) is a rotation \quad (B) doubles area \quad (C) collapses space to a lower dimension \quad (D) preserves length
  \item The eigenvalues of $A$ are the solutions of
    \hfill (A) $A\xx=\zero$ \quad (B) $\det(A-\lambda I)=0$ \quad (C) $A=A^\T$ \quad (D) $A^\T A=I$
  \item The trace of a $2\times2$ matrix equals
    \hfill (A) the product of the eigenvalues \quad (B) the determinant \quad (C) the sum of the eigenvalues \quad (D) the largest eigenvalue
  \item The singular values of a matrix $A$ are
    \hfill (A) the eigenvalues of $A$ \quad (B) $\sqrt{\text{eigenvalues of }A^\T A}$, always $\ge0$ \quad (C) always negative \quad (D) its diagonal entries
  \item $\relu(x)=\max(0,x)$ applied to $x=-4$ gives
    \hfill (A) $-4$ \quad (B) $4$ \quad (C) $0$ \quad (D) $1$
\end{enumerate}

% ======================================================================
\parthead{Part C --- Short Answer \& Computation (48 pts)}
Show your work. ($3$ pts each.)

\begin{enumerate}[leftmargin=1.9em, itemsep=10pt]
  \item \textbf{(U1)} Let $\vv=(5,12)$ and $\ww=(12,-5)$.
        \textbf{(a)} Compute $2\vv+\ww$. \textbf{(b)} Compute $\vv\cdot\ww$ and the lengths $\|\vv\|,\|\ww\|$.
        Are $\vv$ and $\ww$ perpendicular?
        \vspace{2.2cm}
  \item \textbf{(U1)} Let $A=\begin{bsmallmatrix}1&3\\2&6\end{bsmallmatrix}$.
        \textbf{(a)} Compute $A\begin{bsmallmatrix}2\\1\end{bsmallmatrix}$.
        \textbf{(b)} Are the columns independent? What is the rank, and what does the column space look like?
        \vspace{2.2cm}
  \item \textbf{(U2)} Solve by elimination: $\;x+2y=7,\quad 2x+5y=16.$
        \vspace{2.2cm}
  \item \textbf{(U2)} Let $A=\begin{bsmallmatrix}3&1\\5&2\end{bsmallmatrix}$.
        \textbf{(a)} Find $\det A$. \textbf{(b)} Find $A^{-1}$ and check that $A^{-1}A=I$.
        \vspace{2.4cm}
  \item \textbf{(U3)} Reduce $A=\begin{bsmallmatrix}1&3&1\\2&6&3\end{bsmallmatrix}$ and find a special solution of $A\xx=\zero$.
        What is the dimension of the nullspace?
        \vspace{2.4cm}
  \item \textbf{(U3)} The $3\times4$ matrix $A=\begin{bsmallmatrix}1&3&0&2\\0&0&1&4\\0&0&0&0\end{bsmallmatrix}$.
        Give the rank $r$ and the dimensions of all four subspaces: $C(A)$, $N(A)$, the row space $C(A^\T)$, and $N(A^\T)$.
        \vspace{2.6cm}
  \item \textbf{(U4)} Project $\bb=(5,5)$ onto the line through $\avec=(1,3)$. \textbf{(a)} Find the projection $\pp$.
        \textbf{(b)} Find the error $\ee=\bb-\pp$ and verify it is perpendicular to $\avec$.
        \vspace{2.4cm}
  \item \textbf{(U4)} Find the best-fit line $b=C+Dt$ through the points $(0,2),(1,2),(2,8)$ using the normal
        equations $A^\T A\,\hat\xx=A^\T\bb$.
        \vspace{2.8cm}
  \item \textbf{(U5)} Compute the determinants.
        \textbf{(a)} $\det\begin{bsmallmatrix}5&2\\1&3\end{bsmallmatrix}$.
        \textbf{(b)} $\det\begin{bsmallmatrix}2&1&0\\1&3&1\\0&2&1\end{bsmallmatrix}$ (expand along a row).
        \vspace{2.4cm}
  \item \textbf{(U5)} The transformation $A=\begin{bsmallmatrix}3&2\\0&4\end{bsmallmatrix}$.
        \textbf{(a)} What area does a unit square map to? \textbf{(b)} Does $A$ preserve or flip orientation?
        \textbf{(c)} If $\det B=2$, what does $AB$ do to area?
        \vspace{2.4cm}
  \item \textbf{(U6)} Let $A=\begin{bsmallmatrix}3&2\\2&3\end{bsmallmatrix}$. Find its eigenvalues and an eigenvector
        for each. Check your eigenvalues against the trace and determinant.
        \vspace{2.6cm}
  \item \textbf{(U6)} Let $A=\begin{bsmallmatrix}4&2\\1&3\end{bsmallmatrix}$. Use the trace and determinant to find the
        eigenvalues, find an eigenvector for each, and write the diagonalization $A=X\Lambda X^{-1}$ (give $X$ and $\Lambda$).
        \vspace{2.8cm}
  \item \textbf{(U7)} Let $A=\begin{bsmallmatrix}3&1\\-1&-3\end{bsmallmatrix}$. Form $A^\T A$, find its eigenvalues,
        and give the singular values of $A$.
        \vspace{2.6cm}
  \item \textbf{(U7)} \textbf{(a)} Compute the outer product $\uu\vv^\T$ for $\uu=(2,1)$, $\vv=(1,3)$, and say why it is
        rank one. \textbf{(b)} The matrix $A=\begin{bsmallmatrix}6&4\\4&6\end{bsmallmatrix}$ has singular values $10,2$.
        What fraction of the energy does the top rank-one layer keep?
        \vspace{2.4cm}
  \item \textbf{(U8)} Evaluate one neural-network layer $\relu(A\vv+\bb)$ with
        $A=\begin{bsmallmatrix}1&-1\\1&1\end{bsmallmatrix}$, $\vv=\begin{bsmallmatrix}1\\4\end{bsmallmatrix}$,
        $\bb=\begin{bsmallmatrix}2\\-1\end{bsmallmatrix}$. Show $A\vv$, then $A\vv+\bb$, then the ReLU.
        \vspace{2.4cm}
  \item \textbf{(U8)} Four students' two quiz scores are $(1,2),(3,6),(5,4),(7,8)$.
        \textbf{(a)} Find the mean point and the covariance matrix $V$.
        \textbf{(b)} Find the eigenvalues of $V$ and the fraction of the total variance the first principal
        component ($\mathrm{PC}_1$) holds.
        \vspace{2.8cm}
\end{enumerate}

% ======================================================================
\parthead{Part D --- Extended Response (12 pts)}
Explain your reasoning in full sentences. ($2$ pts each.)

\begin{enumerate}[leftmargin=1.9em, itemsep=12pt]
  \item \textbf{(Units 1--3 --- rank)} Explain what the rank $r$ of an $m\times n$ matrix counts (pivots
        $=$ independent columns), and how it fixes the dimensions of the four subspaces. Include why
        $\dim C(A)+\dim N(A)=n$.
        \vspace{2.4cm}
  \item \textbf{(Unit 4 --- least squares)} Explain why the best-fit solution $\hat\xx$ comes from
        \emph{projecting} $\bb$ onto the column space, and why the error must be perpendicular to $C(A)$ ---
        the idea behind the normal equations $A^\T A\,\hat\xx=A^\T\bb$.
        \vspace{2.4cm}
  \item \textbf{(Units 5--6 --- singular)} Explain the chain: $\det A=0$ $\iff$ $A$ is singular $\iff$ the
        columns are dependent $\iff$ $0$ is an eigenvalue. What does $\det A$ measure geometrically?
        \vspace{2.4cm}
  \item \textbf{(Unit 6 --- spectral theorem)} State the spectral theorem for a symmetric matrix $S=S^\T$.
        Why are its eigenvalues real and its eigenvectors perpendicular, and why does this let you write
        $S=Q\Lambda Q^\T$ with an orthogonal $Q$?
        \vspace{2.4cm}
  \item \textbf{(Unit 7 --- the SVD)} Explain why \emph{every} matrix $A$ has a singular value decomposition,
        building the argument from the fact that $A^\T A$ is symmetric (real $\lambda\ge0$, perpendicular
        eigenvectors) so $\sigma=\sqrt\lambda$. Why is this stronger than diagonalization by eigenvectors?
        \vspace{2.4cm}
  \item \textbf{(Unit 8 --- the whole course)} A neural network that learns from data is built entirely from
        tools in this course. Name the linear-algebra idea behind each stage: (i) a layer, (ii) scoring
        predictions with a loss, (iii) choosing the weights, (iv) understanding the shape of the data.
        \vspace{2.6cm}
\end{enumerate}

\end{document}
